\documentclass[10pt]{article}

\usepackage[T1]{fontenc}
\usepackage[utf8]{inputenc}
\usepackage{lmodern}
\usepackage[margin=0.72in]{geometry}
\usepackage{amsmath}
\usepackage{booktabs}
\usepackage{tabularx}
\usepackage{array}
\usepackage{graphicx}
\usepackage{xcolor}
\usepackage{caption}
\usepackage{enumitem}
\usepackage{microtype}
\usepackage{placeins}
\usepackage{hyperref}

\definecolor{accent}{HTML}{245B78}
\hypersetup{colorlinks=true,linkcolor=accent,urlcolor=accent,
  pdfauthor={Tommaso De Santo},
  pdftitle={Property Taxes, Housing Reallocation, and Fertility in E5f}}
\setlength{\parindent}{0pt}
\setlength{\parskip}{0.38em}
\setlist[itemize]{leftmargin=1.3em,itemsep=0.15em,topsep=0.15em}
\setlist[enumerate]{leftmargin=1.45em,itemsep=0.2em,topsep=0.15em}
\captionsetup{font=footnotesize,labelfont=bf,skip=3pt}
\renewcommand{\arraystretch}{1.07}
\newcolumntype{Y}{>{\raggedleft\arraybackslash}X}
\newcommand{\figroot}{../output/model/e5f_post2023_kappa_m5_property_tax_mechanism_decomposition_20260901}
\newcommand{\pp}{\text{ pp}}

\title{\vspace{-2.4em}\textbf{Property Taxes, Housing Reallocation, and Fertility in E5f}\\[-0.05em]
\large Advisor brief}
\author{Tommaso De Santo}
\date{September 2026}

\begin{document}
\maketitle
\vspace{-1.7em}

\section*{The answer}

Doubling the rebated annual property tax from 1 to 2 percent produces a large housing response but only a modest fertility response. Asset prices fall 3--5 percent, ownership falls 2--3 percentage points, and housing services fall 6--8 percent. Annual births per household head rise only 0.2--0.6 percent.

The model's housing--fertility link is present and has the expected sign. The larger problem is that the reform does not deliver more family-sized housing to prospective parents. Some young childless households become owners, but they occupy smaller homes. Young parents lose both ownership and space. The present experiment therefore does not yet reproduce the housing reallocation emphasized by Coven et al.; it cannot tell us how strongly fertility would respond if that reallocation actually reached young constrained families.

\begin{center}
\colorbox{accent!9}{\parbox{0.95\textwidth}{
\textbf{Bottom line.} This is partly calibration and partly theory. The family-housing block is calibrated too weakly at the relevant margins, but the policy also lacks the capitalization, financing, tenure-market, and migration forces needed to generate a strong young-family housing first stage. Positive tenure taste fixes an important numerical problem, but it cannot create those missing mechanisms.
}}
\end{center}

\section*{What was run}

The experiment compares two deterministic perfect-foresight paths beginning from the same reconstructed 2023 distribution. Both span 128 four-year dates and pass the maintained market, fiscal, history, accounting, feasibility, probability, and terminal gates. Post-2023 tenure choice uses $\kappa_{\mathrm{tenure}}=0.0100175$, borrowed from the older M5 calibration. E5f itself was estimated with deterministic tenure choice, so this is a stable conditional sensitivity rather than a jointly recalibrated production result.

\begin{table}[h]
\centering
\caption{Selected effects of a 2 percent rather than 1 percent property tax}
\scriptsize
\begin{tabular}{rrrrrr}
\toprule
Year & Asset price & Ownership & Housing/head & Annual births/head & Equal rebate \\
\midrule
2023 & $-3.196\%$ & $-2.299\pp$ & $-5.525\%$ & $+0.158\%$ & $+82.906\%$ \\
2035 & $-4.574\%$ & $-2.864\pp$ & $-7.864\%$ & $+0.321\%$ & $+75.838\%$ \\
2079 & $-4.344\%$ & $-2.230\pp$ & $-7.572\%$ & $+0.605\%$ & $+76.824\%$ \\
2103 & $-4.224\%$ & $-2.167\pp$ & $-7.515\%$ & $+0.535\%$ & $+77.156\%$ \\
2531 & $-4.164\%$ & $-2.633\pp$ & $-7.538\%$ & $+0.307\%$ & $+77.226\%$ \\
\bottomrule
\end{tabular}
\end{table}

The flow cost of housing rises 17--27 percent even though the purchase price falls. The reform therefore relaxes the down payment modestly but makes housing more expensive to occupy. The equal rebate supplies a broad positive income effect. The small positive birth response is the net of these offsetting forces, not a pure affordability effect.

\newpage
\section*{Housing moves much more than fertility}

\begin{figure}[h]
\centering
\includegraphics[width=0.91\textwidth]{\figroot/paired_transition_effects.pdf}
\caption{Dynamic effects of the property-tax reform}
\end{figure}

The birth response peaks at 0.605 percent in 2079. In levels, annual births per household head rise from 0.0152279 to 0.0153200. The change is 0.0000921 births per head, and $0.0000921/0.0152279=0.00605$. A number near 0.000006 discussed during the solver work was instead a fiscal-equilibrium residual: it measured numerical convergence, not fertility.

The exact state decomposition separates birth behavior for comparable households from movements in the household distribution. If $w_s^j$ is state mass and $p_s^j$ is the four-year birth flow in policy $j$, annual births are $F^j=\frac14\sum_s w_s^jp_s^j$. The exact symmetric decomposition is
\[
F^1-F^0=\frac18\sum_s(w_s^0+w_s^1)(p_s^1-p_s^0)
+\frac18\sum_s(p_s^0+p_s^1)(w_s^1-w_s^0).
\]
The first term is within-state birth behavior; the second is distribution and composition. It is an accounting identity, not a separate causal identification claim.

\begin{table}[h]
\centering
\caption{Exact decomposition of the annual birth-rate response}
\scriptsize
\begin{tabular}{rrrr}
\toprule
Year & Total & Within-state behavior & Distribution/composition \\
\midrule
2023 & $+0.158\%$ & $+0.158\%$ & $-0.000\%$ \\
2035 & $+0.321\%$ & $+0.232\%$ & $+0.089\%$ \\
2051 & $+0.415\%$ & $+0.244\%$ & $+0.170\%$ \\
2079 & $+0.605\%$ & $+0.258\%$ & $+0.347\%$ \\
2103 & $+0.535\%$ & $+0.227\%$ & $+0.308\%$ \\
2531 & $+0.307\%$ & $+0.175\%$ & $+0.132\%$ \\
\bottomrule
\end{tabular}
\end{table}

At the peak, only 0.258 percentage point comes from comparable households changing birth behavior. The remaining 0.347 point---57 percent of the total---comes from changes in the distribution of households across wealth, tenure, age, parity, and child states. The household mechanism is positive, but quantitatively small.

\newpage
\section*{The reform does not expand young-family housing}

\begin{figure}[h]
\centering
\includegraphics[width=0.84\textwidth]{\figroot/fertility_age_decomposition.pdf}
\caption{Age decomposition of the aggregate birth response}
\end{figure}

The fertility response is concentrated at childbearing ages, but the housing incidence is not favorable to families. In 2079, childless households aged 25--34 experience a 2.379 percentage-point ownership increase, yet their housing services fall 7.437 percent. Parents aged 25--34 lose 1.081 ownership points and 6.200 percent of housing; parents aged 35--44 lose 2.943 points and 6.590 percent of housing.

\begin{figure}[h]
\centering
\includegraphics[width=0.82\textwidth]{\figroot/young_household_incidence.pdf}
\caption{Ownership and housing services among young households}
\end{figure}

The reform can move a childless young household from renting to owning while shrinking its housing bundle. It changes tenure without relaxing the family-space margin. No positive-mass young-parent state binds the child-space floor in either policy. No positive-mass at-risk state lies within 0.05 value units of birth indifference; only 2.3--2.6 percent of at-risk mass lies within 0.25 units. There is little household mass at which a modest value change can produce a large birth response.

\newpage
\section*{Why this differs from Coven et al.}

Coven, Golder, Gupta, and Ndiaye study housing allocation across generations; they do not model endogenous fertility. Their mechanism combines lower down payments for young constrained buyers with higher recurring costs that induce older owners to release housing. The age reallocation, rather than the aggregate ownership rate, is the key result.

\begin{table}[h]
\centering
\caption{Housing-allocation mechanisms}
\scriptsize
\begin{tabularx}{\textwidth}{>{\raggedright\arraybackslash}p{0.22\textwidth}XX}
\toprule
Object & Coven et al. & Current E5f sensitivity \\
\midrule
Property-tax change & 0.8 to 2.0 percent & 1.0 to 2.0 percent \\
Supply elasticity & 0.232 & 1.75 diagnostic value \\
Asset-price response & $-12.6$ percent & $-3.2$ percent initially; $-4.7$ percent at peak \\
Payment-to-income limit & 36 percent and active & Inactive \\
Owner and rental housing & Segmented stocks & One housing-service market \\
Spatial adjustment & Two-region migration extension & Absent from this run \\
Young incidence & Ownership rises at ages 25--34 while old ownership falls sharply & Childless young eventually own more, but all young groups consume less housing and parents own less \\
\bottomrule
\end{tabularx}
\end{table}

The differences locate the missing first stage. More elastic supply produces less price capitalization and hence less down-payment relief. Without an active payment-to-income constraint, the higher tax does not sharply induce cash-flow-constrained older owners to downsize. A common owner--renter housing-service market makes tenure less consequential for space. Without migration, households cannot sort toward locations where the tax releases larger homes.

\section*{Why the fertility block amplifies so little}

Fertility depends on the difference between the optimized value of attempting a birth and waiting. Housing prices affect both branches. If the branches choose similar interior housing bundles, their common housing-cost exposure largely cancels from the birth value gap. Housing becomes a strong fertility margin when a child creates a large additional need for space, one branch is constrained, or tenure supplies a child-specific service.

The calibration currently supplies little of that amplification. The first-child housing floor rises by only $0.19945+0.31527=0.51471$ room-equivalent units, while young parents consume roughly 5 units as renters and 8--9 as owners. The model predicts only a 0.380-room increase after a first birth against 0.720 in the data, while predicting 6.710 average rooms against 5.780. Prospective parents begin with too much space and add too little after childbirth.

The target set also contains no moment measuring fertility responses to plausibly exogenous rents, prices, taxes, supply, or down-payment relief. ``Housing after a birth'' and ``births after a housing-cost shock'' are different derivatives. The latter is presently model-implied rather than directly identified.

\section*{Recommended sequence}

\begin{enumerate}
  \item Establish the Coven-style first stage with bounded sensitivities for supply elasticity, an explicitly calibrated payment-to-income constraint, and owner--renter segmentation. Report age-specific housing services, not only ownership.
  \item Repair the first-birth room response and average-room fit.
  \item Add a moment identifying fertility responses to exogenous housing-cost or financing variation.
  \item Re-estimate the positive tenure taste jointly with E5f.
  \item Add child-specific tenure services only if young parents receive more housing and fertility still barely responds.
\end{enumerate}

The positive tenure taste is worth retaining because it smooths an economically relevant discrete choice and stabilizes the transition. But it cannot by itself add the financial, spatial, or family-housing mechanisms above.

\newpage
\section*{Compact calibration appendix}

The active E5f transition loss is 36.0992. The two housing-size rows contribute 26.304 points, making the family-housing fit the dominant calibration tension.

\begin{table}[h]
\centering
\caption{Complete active target fit}
\scriptsize
\begin{tabularx}{\textwidth}{>{\raggedright\arraybackslash}Xrrrrrr}
\toprule
Moment & Target & Model & Gap & Weight & Std. gap & Contribution \\
\midrule
Total fertility rate & 1.9180 & 1.9183 & $+0.0003$ & 1,425.74 & $+0.009$ & 0.000 \\
Childless rate & 0.1880 & 0.1859 & $-0.0021$ & 17,180.74 & $-0.280$ & 0.078 \\
Mean age at first birth & 26.0446 & 26.3273 & $+0.2827$ & 44.44 & $+1.885$ & 3.552 \\
First births at age 30+ & 0.2603 & 0.2421 & $-0.0182$ & 10,000.00 & $-1.821$ & 3.317 \\
Rooms gained after first birth & 0.7202 & 0.3798 & $-0.3405$ & 137.57 & $-3.993$ & 15.946 \\
Rooms: parity 3+ minus parity 1--2 & 0.3677 & 0.3593 & $-0.0084$ & 2,958.51 & $-0.459$ & 0.211 \\
Family ownership gap & 0.1677 & 0.1769 & $+0.0092$ & 14,229.59 & $+1.098$ & 1.206 \\
Ownership rate & 0.5755 & 0.5461 & $-0.0294$ & 1,207.85 & $-1.022$ & 1.045 \\
Mean occupied rooms, ages 18--85 & 5.7800 & 6.7101 & $+0.9301$ & 11.97 & $+3.218$ & 10.359 \\
Wealth / annual gross earnings & 6.8731 & 6.7664 & $-0.1067$ & 6.29 & $-0.268$ & 0.072 \\
Annual bequest flow / wealth & 0.0088 & 0.0086 & $-0.0002$ & 5,165,289.26 & $-0.402$ & 0.162 \\
Old-age wealth p90/p50, ages 76--84 & 3.4481 & 3.3963 & $-0.0518$ & 56.96 & $-0.391$ & 0.153 \\
\bottomrule
\end{tabularx}
\end{table}

\begin{table}[h]
\centering
\caption{Complete active free-parameter table}
\scriptsize
\begin{tabular}{lrrrrl}
\toprule
Parameter & Estimate & Lower & Upper & Transform & Near bound \\
\midrule
$\beta_{\mathrm{annual}}$ & 0.994959 & 0.94 & 0.9995 & discount & No \\
$\kappa_{\mathrm{fert}}$ & 2.267631 & 0.02 & 50.0 & log & No \\
$\kappa_{\mathrm{fert,cont}}$ & 1.765427 & 0.02 & 50.0 & log & No \\
$\chi$ & 1.033945 & 0.10 & 5.0 & log & No \\
$H_0$ & 16.381189 & 0.20 & 80.0 & log & No \\
$\theta_0$ & 0.563219 & 0.00 & 8.0 & soft zero & No \\
$\theta_1$ & 0.099707 & 0.02 & 16.0 & log & Yes \\
$\bar h_{\mathrm{child}}$ & 0.315267 & 0.10 & 1.8 & log & No \\
First-birth fixed cost & 4.545167 & 0.00 & 8.0 & soft zero & No \\
First-child housing jump & 0.199446 & 0.00 & 0.5 & soft zero & No \\
$\Delta\psi_{\mathrm{child},2023}$ & $-0.350000$ & $-1.5$ & 0.2 & asinh & No \\
\bottomrule
\end{tabular}
\end{table}

\vfill
\footnotesize
\textit{Source and verification.} All paths, state decompositions, tables, and figures come from the hash-audited packet \path{output/model/e5f_post2023_kappa_m5_property_tax_mechanism_decomposition_20260901/}. The collector verifies equilibrium, accounting, feasibility, probability, state-grid equality, and exact decomposition add-up. Coven et al.: Joshua Coven, Sebastian Golder, Arpit Gupta, and Abdoulaye Ndiaye, ``Property Taxes and Housing Allocation Under Financial Constraints,'' August 1, 2026, \url{https://abdouecon.github.io/research/papers/Property_Tax.pdf}.

\end{document}
